\documentclass[11pt]{article}
\usepackage[margin=1in]{geometry}
\usepackage{amsmath,amssymb,amsthm}
\usepackage{xurl}
\usepackage{hyperref}
\sloppy
\let\oldus\_ \renewcommand{\_}{\oldus\allowbreak}
\newtheorem{theorem}{Theorem}
\newtheorem{lemma}[theorem]{Lemma}
\newtheorem{corollary}[theorem]{Corollary}
\theoremstyle{definition}
\newtheorem{definition}[theorem]{Definition}
\newtheorem{remark}[theorem]{Remark}
\title{Conjecture 00000007864 is false:\\ the men-optimal stable matching is exactly uniform,\\ so the entropy deficit constant is $1$, not $\pi^2/6$}
\author{Jackmeson1}
\date{}
\begin{document}
\maketitle

\section{The conjecture and the reading}

\textbf{Statement (English text of the conjecture).} ``The spousal distribution concentration of random stable
matching: the entropy $H(\sigma_m)$ of the men-optimal matching under uniform $n$ by $n$ preferences (the
Gale-Shapley expected stable outcome). Conjecture: $H(\sigma_m) = n\log n - \Theta(n)$ (a linear entropy deficit)
with deficit constant $c^* = \pi^2/6$, given by the eigenvalue distribution of a Hankel matrix and a Kotochyk-type
integral; and the longest increasing subsequence of $\sigma_m$ has length $2\sqrt n + O(n^{1/6})$ with
Tracy-Widom fluctuations, [...]'' The Chinese text states the same: the deficit constant is $c^*=\pi^2/6$.

\textbf{Reading.} There are $n$ men and $n$ women, both indexed by $\{0,\dots,n-1\}$. Every person has a strict
preference list over the other side; all $2n$ lists are independent and uniformly distributed over the $n!$
orderings, i.e.\ the profile is uniform over all $(n!)^{2n}$ profiles. $\sigma_m$ is the men-optimal stable
matching, $H(\sigma_m)$ is the Shannon entropy (natural logarithm) of the law of $\sigma_m$, and the deficit is
$D_n = n\log n - H(\sigma_m)$. The refuted clause is ``deficit constant $c^*=\pi^2/6$'', read as
$D_n/n \to \pi^2/6$. We prove $D_n/n \le 1 < \pi^2/6$ for every $n\ge 1$ and $D_n/n\to 1$; so the limit, limsup,
liminf and asymptotic-equivalence ($D_n\sim c^* n$) readings of the deficit constant all fail. With base-$2$
logarithms throughout, $D_n/n \to 1/\ln 2\approx 1.4427$, again not $\pi^2/6$.

\textbf{What is not refuted.} The linear-deficit statement $D_n=\Theta(n)$ is true (it follows from our bounds).
The clause on the longest increasing subsequence is not addressed. The conjecture is refuted as a conjunction,
through its clause $c^*=\pi^2/6$.

\textbf{Conventions.} A preference list is encoded as a rank bijection, rank $0$ being the most preferred; rank
bijections correspond one-to-one to strict total orders, so the uniform law on rank bijections is the uniform law
on preference lists. Matchings are bijections from men to women. $\log$ is the natural logarithm and
$0\log 0=0$ (Mathlib's \texttt{Real.negMulLog}).

\section{Definitions (as in the Lean file)}

\begin{definition}
A \emph{profile} $P=(r,s)$ consists of rank bijections $r_m:\{\text{women}\}\to\{0,\dots,n-1\}$ for each man $m$
and $s_w:\{\text{men}\}\to\{0,\dots,n-1\}$ for each woman $w$ (\texttt{Profile n}). A \emph{matching} is a
permutation $\mu$: man $m$ is married to woman $\mu(m)$ (\texttt{Matching n}). The pair $(m,w)$ \emph{blocks}
$\mu$ if $r_m(w)<r_m(\mu(m))$ and $s_w(m)<s_w(\mu^{-1}(w))$ (\texttt{Blocks}); $\mu$ is \emph{stable} if no pair
blocks it (\texttt{IsStable}); $\mu$ is \emph{men-optimal} if it is stable and $r_m(\mu(m))\le r_m(\nu(m))$ for
every stable $\nu$ and every man $m$ (\texttt{IsMenOptimal}).
\end{definition}

This is the textbook notion: a matching is stable ``when there does not exist any pair (A, B) where both prefer
each other to their current partner under the matching'' [W].

\begin{definition}
The law of $\sigma_m$ is the push-forward of the uniform distribution on profiles under the men-optimal map:
\texttt{lawMenOpt n = (PMF.uniformOfFintype (Profile n)).map menOpt}. The entropy of a distribution $p$ on a
finite set is $H(p)=\sum_a -p(a)\log p(a)$ (\texttt{entropy}); \texttt{deficit n} $= n\log n - H(\texttt{lawMenOpt n})$.
\end{definition}

\section{Gale-Shapley: existence and uniqueness}

Gale and Shapley showed that stable matchings always exist, and the men-proposing algorithm ``always yields the
one that is best for all men among all stable matchings'' [W]. We re-prove this in Lean.

\begin{theorem}[\texttt{exists\_menOptimal}, \texttt{menOptimal\_unique}]
Every profile has exactly one men-optimal stable matching.
\end{theorem}

\begin{proof}
For a set $R$ of pairs (``$w$ has rejected $m$''), say $m$ \emph{proposes} to $w$ if $(m,w)\notin R$ and every
woman $m$ ranks above $w$ has rejected him. The invariant $\mathrm{Inv}(R)$ says: (1) no pair of $R$ is married in
any stable matching; (2) if $(m,w)\in R$ then some $m'$ proposes to $w$ and $s_w(m')<s_w(m)$. The empty set
satisfies it. Take $R$ satisfying $\mathrm{Inv}$ of maximal size. Every man proposes to someone: otherwise every
woman holds a proposal from a man other than $m$, and these men are distinct (a man proposes to one woman), so
$n$ women would be covered by $n-1$ men. Let $p(m)$ be $m$'s proposal. If $p$ is injective it is a matching. It is
stable: if $(m,w)$ blocked it, $w$ would have rejected $m$, so by (2) she holds a proposal $p(m')=w$ with
$s_w(m')<s_w(m)$, i.e.\ she prefers her husband. It is men-optimal: if $r_m(\nu(m))<r_m(p(m))$ then
$(m,\nu(m))\in R$, so $\nu$ is not stable by (1). If $p$ is not injective, two men $l\ne b$ propose to the same $w$,
and $w$ prefers $b$. Then $R\cup\{(l,w)\}$ still satisfies $\mathrm{Inv}$: for (1), if a stable $\nu$ had
$\nu(l)=w$, then $\nu(b)\ne w$, $(b,\nu(b))\notin R$ by (1), so $b$ ranks $w$ strictly above $\nu(b)$, and $w$
prefers $b$ to $l=\nu^{-1}(w)$: $(b,w)$ blocks $\nu$. For (2), every proposal other than $l$'s is unchanged, and
$b$ witnesses the new pair and any old pair whose witness was $l$ (transitivity). This contradicts maximality.
Uniqueness: two men-optimal matchings give every man partners of equal rank, hence the same partner.
\end{proof}

\section{Exact uniformity of $\sigma_m$}

For a permutation $\tau$ of the women let $\tau\cdot P$ rename woman $w$ as $\tau(w)$: man $m$ ranks $\tau(w)$ as he
ranked $w$, and woman $\tau(w)$ gets $w$'s list (\texttt{relabel}).

\begin{lemma}[\texttt{blocks\_relabel}, \texttt{stable\_relabel}, \texttt{menOpt\_relabel}]
$(m,w)$ blocks $\tau\circ\mu$ under $\tau\cdot P$ iff $(m,\tau^{-1}w)$ blocks $\mu$ under $P$. Hence $\tau\circ\mu$
is stable for $\tau\cdot P$ iff $\mu$ is stable for $P$, and $\sigma_m(\tau\cdot P)=\tau\circ\sigma_m(P)$.
\end{lemma}

\begin{theorem}[\texttt{law\_menOpt\_eq\_uniform}]
Under uniform independent preferences, $\sigma_m$ is uniformly distributed on $S_n$.
\end{theorem}

\begin{proof}
$P\mapsto\tau\cdot P$ is a bijection of the profiles with inverse $P\mapsto\tau^{-1}\cdot P$
(\texttt{relabel\_relabel}), and it maps the fibre $\{P:\sigma_m(P)=\sigma\}$ onto the fibre of $\tau\circ\sigma$
(\texttt{card\_fiber\_trans}). Taking $\tau=\sigma'\circ\sigma^{-1}$, all $n!$ fibres have the same size, so each has
$(n!)^{2n}/n!$ elements (\texttt{card\_fiber\_mul}) and probability $1/n!$ (\texttt{lawMenOpt\_toReal}).
\end{proof}

\section{Entropy and the deficit}

\begin{theorem}[\texttt{entropy\_menOpt}, \texttt{deficit\_eq}]
$H(\sigma_m)=\log n!$ and $D_n=n\log n-\log n!$.
\end{theorem}

\begin{theorem}[\texttt{deficit\_le}, \texttt{deficit\_ge}]
$D_n\le n$ for all $n$, and $D_n\ge n-1-\tfrac12\log n$ for $n\ge 1$.
\end{theorem}

\begin{proof}
Upper bound: $n^n/n!\le e^n$ (Mathlib \texttt{Real.pow\_div\_factorial\_le\_exp}); take logarithms. Lower bound:
Mathlib's Stirling sequence $s_k=k!/(\sqrt{2k}(k/e)^k)$ is antitone for $k\ge1$ (\texttt{Stirling.stirlingSeq'\_antitone}),
so $s_n\le s_1=e/\sqrt2$. With \texttt{Stirling.log\_stirlingSeq\_formula},
$\log n! = \log s_n+\tfrac12\log(2n)+n\log n-n\le 1+\tfrac12\log n+n\log n-n$.
\end{proof}

\begin{theorem}[\texttt{conjecture\_7864\_false}]
For every $n\ge1$, $D_n/n\le 1<\pi^2/6$; moreover $D_n/n\to1$, and $D_n/n$ does not tend to $\pi^2/6$.
\end{theorem}

\begin{proof}
The first claim is \texttt{deficit\_le} with $\pi>3$. The squeeze $1-\frac1n-\frac{\log n}{2n}\le D_n/n\le1$ and
$\log n/n\to0$ give $D_n/n\to1$ (\texttt{deficit\_div\_tendsto}); limits are unique and $1\ne\pi^2/6$.
\end{proof}

\begin{corollary}[\texttt{bits\_version}]
With base-$2$ logarithms in both $n\log n$ and the entropy, the normalised deficit tends to $1/\ln 2$, which is not
$\pi^2/6$ (using $\ln 2>0.6931$ and $\pi>3.14$).
\end{corollary}

\begin{remark}
(Not formalised.) If $H(\sigma_m)$ were read as the sum over men of the entropies of their wives' marginal laws,
uniformity gives marginals uniform on $n$ women, so this sum is $n\log n$ and the deficit is $0$, which is not
$\Theta(n)$ either.
\end{remark}

\section{Lean formalisation and axiom audit}

File \texttt{Conjecture7864/Basic.lean} (390 lines, only import \texttt{Mathlib}, Lean 4.33.1 with the Mathlib
\texttt{v4.33.1} tag). Main theorem:
\begin{verbatim}
theorem conjecture_7864_false :
    (forall n : Nat, 1 <= n ->
        deficit n / n <= 1 /\ deficit n / n < Real.pi ^ 2 / 6) /\
    Tendsto (fun n : Nat => deficit n / n) atTop (nhds 1) /\
    not Tendsto (fun n : Nat => deficit n / n) atTop (nhds (Real.pi ^ 2 / 6))
\end{verbatim}
(shown in ASCII; the file uses the usual Unicode symbols). \texttt{\#print axioms} for
\texttt{conjecture\_7864\_false}, \texttt{bits\_version} and \texttt{law\_menOpt\_eq\_uniform} reports only
\texttt{propext}, \texttt{Classical.choice}, \texttt{Quot.sound}. There is no \texttt{sorry} and no
\texttt{native\_decide}.

\section*{Source}
[W] Wikipedia, ``Stable matching problem'', \url{https://en.wikipedia.org/wiki/Stable_matching_problem}. Quoted:
``In 1962, David Gale and Lloyd Shapley proved that, for any equal number in different groups, in the context of
college admissions and individuals wanting marriage it is always possible to solve as matched couples to make all
resultant pairings / matched factors stable.'' ``Among all possible different stable matchings, it always yields
the one that is best for all men among all stable matchings, and worst for all women.'' ``In other words, a
matching is stable when there does not exist any pair (A, B) where both prefer each other to their current partner
under the matching.''

\end{document}
